(1) 写 \(E=-V_0+\hbar^2k^2/(2m)=-\hbar^2\kappa^2/(2m)\)，\(k,\kappa>0\) 且 \(k^2+\kappa^2=2mV_0/\hbar^2\)。偶态内部 \(A\cos kx\)，奇态 \(A\sin kx\)，外侧与 \(e^{-\kappa(|x|-a)}\) 连接。导数连续得
\[
k\tan(ka)=\kappa\quad(\text{偶}),\qquad
-k\cot(ka)=\kappa\quad(\text{奇}).
\]
(2) 设 \(z=ka,z_0=a\sqrt{2mV_0}/\hbar,\ y=\kappa a=\sqrt{z_0^2-z^2}\)，故
\(z\tan z=y\) 或 \(-z\cot z=y\)，\(0<z<z_0\)。每个tan/cot允许正右侧的连续半周期区间至多有一交点，交点按偶、奇、偶、奇交替。任意 \(z_0>0\) 第一偶态始终存在。
(3) 当新态从 \(E\uparrow0\) 进入束缚谱时 \(y=0,z=z_0\)。偶阈值满足 \(\tan z_0=0\)，故除零外 \(z_0=\pi,2\pi,\dots\)；奇阈值满足 \(\cot z_0=0\)，故 \(z_0=\pi/2,3\pi/2,\dots\)。合并为 \(z_0=n\pi/2\)（正整数n）。在恰为阈值处 \(\kappa=0\)，外部波函数不衰减，不能平方可积，故真正负能束缚态数为
\[
N_\mathrm{bound}=\left\lceil {2z_0\over\pi}\right\rceil
\]
（当 \(z_0\) 恰为整数个 \(\pi/2\) 时，阈值新态不计，旧的数目恰为该整数；否则为上一整数加一）。
(4) \(z_0\to\infty\) 时穿透长度 \(1/\kappa\to0\) 对固定低模成立，边界近似Dirichlet，内部 \(k_n\to n\pi/(2a)\)（\(n=1,2,\dots\)），故
\[
E_n+V_0\to{\hbar^2\pi^2n^2\over8ma^2}.
\]
浅阱 \(z_0\ll1\) 时基态偶方程 \(z\tan z\simeq z^2\simeq z_0\!^2\) 给 \(y\simeq z_0^2\)，所以 \(E_0=-\hbar^2\kappa^2/(2m)\simeq-2mV_0^2a^2/\hbar^2\)。
